\RequirePackage{fix-cm}
\documentclass[17pt,a4paper,oneside]{extarticle}

\usepackage[a4paper,top=2.2cm,bottom=2.2cm,left=2.7cm,right=2.2cm]{geometry}
\usepackage{fontspec}
\usepackage{polyglossia}
\setdefaultlanguage{thai}
\setotherlanguage{english}
\XeTeXlinebreaklocale "th"
\XeTeXlinebreakskip=0pt plus 1pt
\defaultfontfeatures{Ligatures=TeX}
\setmainfont[
  Path=fonts/,
  UprightFont=THSarabunNew.ttf,
  BoldFont=THSarabunNew Bold.ttf,
  ItalicFont=THSarabunNew Italic.ttf,
  BoldItalicFont=THSarabunNew BoldItalic.ttf
]{THSarabunNew}
\newfontfamily\thaifont[
  Path=fonts/,
  UprightFont=THSarabunNew.ttf,
  BoldFont=THSarabunNew Bold.ttf,
  ItalicFont=THSarabunNew Italic.ttf,
  BoldItalicFont=THSarabunNew BoldItalic.ttf
]{THSarabunNew}
\newfontfamily\englishfont[
  Path=fonts/,
  UprightFont=THSarabunNew.ttf,
  BoldFont=THSarabunNew Bold.ttf,
  ItalicFont=THSarabunNew Italic.ttf,
  BoldItalicFont=THSarabunNew BoldItalic.ttf
]{THSarabunNew}

\usepackage{setspace}
\usepackage{enumitem}
\usepackage[indentafter]{titlesec}
\usepackage{array}
\usepackage{tabularx}
\usepackage{longtable}
\usepackage[table]{xcolor}
\usepackage{fancyhdr}
\usepackage{ragged2e}

% ข้อมูลสำหรับแบบเสนอหัวข้องานวิจัย CS1
% แก้เฉพาะค่าด้านขวาของแต่ละคำสั่งก่อนส่งเอกสารจริง

\newcommand{\proposalcode}{CS1}
\newcommand{\proposalstudentid}{6611501658}
\newcommand{\proposalstudentname}{ทิวัตถ์ เลิศชัย}
\newcommand{\proposalstudyprogram}{ภาคในเวลา}
\newcommand{\proposalyearlevel}{4}
\newcommand{\proposalgroup}{วท.บ.661(4)/1}

\newcommand{\proposaltitleth}{การพัฒนาระบบจัดการและบันทึกการฝึกงานแบบมีส่วนร่วมหลายฝ่ายสำหรับคณะวิทยาศาสตร์ มหาวิทยาลัยราชภัฏจันทรเกษม}
\newcommand{\proposaltitleen}{Development of a Multi-stakeholder Internship Management and Logbook System for the Faculty of Science, Chandrakasem Rajabhat University}

\newcommand{\proposaladvisor}{ผู้ช่วยศาสตราจารย์วรรณา วิโรจน์แดนไทย}
\newcommand{\proposalchair}{รอตรวจสอบ}

% ช่วงเวลาด้านล่างเป็นแผนเบื้องต้น ให้แก้เป็นเดือนจริงก่อนส่ง
\newcommand{\proposalperiod}{8 เดือน นับจากวันที่หัวข้อได้รับอนุมัติ}

% เทคโนโลยีของระบบจริงยังไม่ปรากฏในเอกสารต้นทาง จึงตั้งค่าแบบไม่ผูกกับผลิตภัณฑ์
\newcommand{\proposalhardware}{คอมพิวเตอร์สำหรับพัฒนาและทดสอบระบบที่มีหน่วยประมวลผลแบบ 64 บิต หน่วยความจำอย่างน้อย 8 GB พื้นที่ว่างอย่างน้อย 20 GB อุปกรณ์พกพาที่มีเว็บเบราว์เซอร์และระบบระบุตำแหน่งสำหรับทดสอบการลงเวลา และเครื่องแม่ข่ายหรือบริการคลาวด์สำหรับติดตั้งระบบทดสอบ}
\newcommand{\proposalsoftware}{ระบบปฏิบัติการที่รองรับเครื่องมือพัฒนา เว็บเบราว์เซอร์รุ่นปัจจุบัน Git โปรแกรมแก้ไขซอร์สโค้ด Docker และ Docker Compose รวมถึงฐานข้อมูล เฟรมเวิร์กส่วนหน้า และเฟรมเวิร์กฝั่งเซิร์ฟเวอร์ตามรุ่นจริงของ Trainy ซึ่งต้องระบุชื่อและหมายเลขรุ่นก่อนส่งข้อเสนอ}


\setstretch{1.0}
% เพิ่มความยืดหยุ่นในการตัดบรรทัดภาษาไทย เพื่อไม่ให้ข้อความล้นขอบกระดาษ
\setlength{\emergencystretch}{4em}
\sloppy
\setlength{\parindent}{1.27cm}
\setlength{\parskip}{0pt}
\setlist{nosep,leftmargin=1.15cm,labelsep=.35cm}
\renewcommand{\arraystretch}{1.25}
\setlength{\headheight}{22pt}
% ใช้ลำดับเลขหัวข้อแบบเดียวกับเล่มวิจัย และรองรับถึง 2.1.1.1
\setcounter{secnumdepth}{4}
\titleformat{\section}
  {\bfseries\fontsize{16pt}{16pt}\selectfont}
  {\thesection.}{0.45em}{}
\titleformat{\subsection}
  {\normalfont\fontsize{16pt}{16pt}\selectfont}
  {\thesubsection}{0.45em}{}
\titleformat{\subsubsection}
  {\normalfont\fontsize{16pt}{16pt}\selectfont}
  {\thesubsubsection}{0.45em}{}
\titleformat{\paragraph}[block]
  {\normalfont\fontsize{16pt}{16pt}\selectfont}
  {\theparagraph}{0.45em}{}
\titlespacing*{\section}{0pt}{12pt}{0pt}
\titlespacing*{\subsection}{0.95cm}{0pt}{0pt}
\titlespacing*{\subsubsection}{2.22cm}{0pt}{0pt}
\titlespacing*{\paragraph}{3.49cm}{0pt}{0pt}
\pagestyle{fancy}
\fancyhf{}
\fancyhead[R]{\proposalcode}
\fancyfoot[C]{\thepage}
\renewcommand{\headrulewidth}{0pt}
\renewcommand{\footrulewidth}{0pt}

\newcommand{\fieldline}[1]{\textbf{#1}}
\newcommand{\pending}[1]{\textcolor{red!70!black}{#1}}
\newcommand{\studentvalue}[1]{#1}
\newcommand{\studentpair}[2]{\textbf{#1}\hspace{0.18cm}\studentvalue{#2}}
% รูปแบบช่องลงนามจาก front-matter/approval-page.tex
\newcommand{\signatureblock}[2]{%
  \noindent\begin{tabularx}{\textwidth}{@{}X>{\centering\arraybackslash}p{7cm}@{\hspace{0.5em}}p{3.3cm}@{}}
    & \hrulefill & #1 \\
    & (#2) & \\
  \end{tabularx}}

\begin{document}

\begin{center}
  {\bfseries\fontsize{16pt}{19.2pt}\selectfont แบบเสนอหัวข้องานวิจัยทางวิทยาการคอมพิวเตอร์และปัญญาประดิษฐ์\par}
  \vspace{2pt}
  {\bfseries คณะวิทยาศาสตร์ มหาวิทยาลัยราชภัฏจันทรเกษม\par}
\end{center}

\section{ข้อมูลนักศึกษา}
\noindent
\makebox[6.6cm][l]{\studentpair{รหัสนักศึกษา}{\proposalstudentid}}%
\studentpair{ชื่อ--นามสกุล}{\proposalstudentname}\\[4pt]
\noindent
\makebox[3.8cm][l]{\studentpair{นักศึกษา}{\proposalstudyprogram}}%
\makebox[2.8cm][l]{\studentpair{ชั้นปี}{\proposalyearlevel}}%
\studentpair{หมู่เรียน}{\proposalgroup}

\section{ชื่อหัวข้อที่นำเสนอ}
\noindent\begin{tabularx}{\textwidth}{@{}>{\bfseries}p{2.4cm}@{\hspace{0.25cm}}>{\RaggedRight\hyphenpenalty=10000\exhyphenpenalty=10000\arraybackslash}X@{}}
ภาษาไทย & \proposaltitleth \\
ภาษาอังกฤษ & \textenglish{\proposaltitleen}
\end{tabularx}

\section{ความเป็นมาและความสำคัญของปัญหา}
การฝึกงานและการศึกษาเชิงบูรณาการกับการทำงานเชื่อมโยงการเรียนในมหาวิทยาลัยกับการปฏิบัติงานจริง นักศึกษาต้องประสานงานกับอาจารย์ เจ้าหน้าที่ และสถานประกอบการ ตั้งแต่การสมัครและอนุมัติสถานที่ฝึกงาน การบันทึกเวลาและกิจกรรม การติดตามความก้าวหน้า การนิเทศ จนถึงการประเมินและสรุปผล จึงเป็นกระบวนการหลายฝ่ายที่ต้องใช้ข้อมูลร่วมกันอย่างต่อเนื่อง

เมื่อข้อมูลกระจายอยู่ในเอกสาร แบบฟอร์ม และช่องทางติดต่อหลายระบบ อาจเกิดข้อมูลซ้ำซ้อน ตรวจสอบสถานะได้ยาก และติดตามการปฏิบัติงานไม่ต่อเนื่อง ระบบสารสนเทศด้านการฝึกงานจึงควรเป็นศูนย์กลางของข้อมูล เอกสาร การอนุมัติ การติดตาม และการประเมิน พร้อมกำหนดสิทธิ์ให้เหมาะกับบทบาทของผู้ใช้

Trainy เป็นระบบที่พัฒนาขึ้นเพื่อสนับสนุนกระบวนการดังกล่าว โดยรองรับนักศึกษา อาจารย์ \mbox{ผู้ประสานงาน} เจ้าหน้าที่มหาวิทยาลัย ผู้ดูแลบริษัท และพี่เลี้ยงในสถานประกอบการ อย่างไรก็ตาม ระบบรายงานความก้าวหน้ารุ่นปัจจุบันยังเก็บข้อความสรุปและจำนวนชั่วโมงเป็นหลัก จึงยังไม่ครอบคลุมรายละเอียดของสมุดบันทึกฝึกงานแบบมีโครงสร้าง เช่น กิจกรรม สิ่งที่เรียนรู้ ปัญหา วิธีแก้ไข ทักษะ และผลลัพธ์ของงาน อีกทั้งยังต้องประเมินว่าผู้ใช้แต่ละบทบาทสามารถทำภารกิจหลักได้จริงหรือไม่

งานวิจัยนี้จึงมุ่งศึกษา ออกแบบ และพัฒนา Trainy สำหรับคณะวิทยาศาสตร์ มหาวิทยาลัยราชภัฏจันทรเกษม พร้อมประเมินการทำภารกิจและความสามารถในการใช้งานกับผู้เกี่ยวข้อง 4 กลุ่ม ได้แก่ นักศึกษา อาจารย์ เจ้าหน้าที่ และผู้ประกอบการหรือพี่เลี้ยง ผลการศึกษาจะใช้เป็นหลักฐานสำหรับปรับปรุงระบบก่อนนำไปใช้จริงในวงกว้าง โดยไม่อ้างอิงผลแทนมหาวิทยาลัยอื่น

\section{วัตถุประสงค์}
\begin{enumerate}
  \item เพื่อศึกษากระบวนการ ความต้องการ และปัญหาในการจัดการและบันทึกการฝึกงานของคณะวิทยาศาสตร์ มหาวิทยาลัยราชภัฏจันทรเกษม
  \item เพื่อออกแบบและพัฒนา Trainy ให้รองรับกระบวนการฝึกงานตั้งแต่เริ่มต้นจนถึงการประเมินและสรุปผล โดยเชื่อมโยงนักศึกษา อาจารย์ เจ้าหน้าที่ และสถานประกอบการ
  \item เพื่อประเมินประสิทธิผลในการทำภารกิจ ปัญหาการใช้งาน และความสามารถในการใช้งานของ Trainy จากผู้ใช้ที่เกี่ยวข้อง
  \item เพื่อจัดทำข้อเสนอแนะสำหรับปรับปรุงระบบบันทึกและติดตามการฝึกงานให้เหมาะสมกับการใช้งานจริง
\end{enumerate}

\section{หลักการ ทฤษฎี และเหตุผล}
งานวิจัยใช้แนวคิดระบบสารสนเทศเพื่อรวมข้อมูลและกระบวนการของผู้เกี่ยวข้องหลายฝ่าย แนวคิดการออกแบบโดยยึดผู้ใช้เป็นศูนย์กลางเพื่อศึกษาภารกิจและความต้องการตามบทบาท และแนวคิดสมุดบันทึกดิจิทัลเพื่อจัดเก็บกิจกรรม การเรียนรู้ ปัญหา วิธีแก้ไข และข้อเสนอแนะอย่างตรวจสอบย้อนหลังได้ การพัฒนาจะดำเนินในลักษณะวิจัยและพัฒนาแบบกรณีศึกษา โดยใช้ข้อมูลเชิงปริมาณร่วมกับข้อมูลเชิงคุณภาพ

การประเมินจะพิจารณาความสำเร็จของภารกิจ เวลา ข้อผิดพลาด และการขอความช่วยเหลือ ร่วมกับแบบประเมิน System Usability Scale (SUS) จำนวน 10 ข้อและคำถามปลายเปิด วิธีดังกล่าวเหมาะกว่าการใช้คะแนนความพึงพอใจเพียงค่าเดียว เพราะช่วยระบุทั้งระดับการใช้งานโดยรวมและจุดติดขัดในขั้นตอนจริง ทั้งนี้จะรายงานผลแยกตามบทบาทเพื่อไม่ให้จำนวนผู้ใช้กลุ่มใหญ่กลบปัญหาของกลุ่มย่อย

\section{ระยะเวลาดำเนินการ}
กำหนดระยะเวลาดำเนินการเบื้องต้น \fieldline{\proposalperiod} โดยนับเดือนที่ 1 เป็นเดือนแรกหลังหัวข้อได้รับอนุมัติ กำหนดการจริงต้องปรับให้ตรงกับปฏิทินการศึกษาและช่วงฝึกงานของกลุ่มตัวอย่าง

\clearpage
\begin{center}
\small
\textbf{แผนการดำเนินงานรายเดือน (Gantt chart)}\\[3pt]
\begin{tabularx}{\textwidth}{|>{\RaggedRight\arraybackslash}X|*{8}{>{\centering\arraybackslash}p{0.65cm}|}}
\hline
\rowcolor{gray!20}\textbf{กิจกรรม} & \textbf{1} & \textbf{2} & \textbf{3} & \textbf{4} & \textbf{5} & \textbf{6} & \textbf{7} & \textbf{8} \\
\hline
ศึกษากระบวนการ เอกสาร และผู้เกี่ยวข้อง & \cellcolor{gray!45} & \cellcolor{gray!45} & & & & & & \\
\hline
วิเคราะห์ความต้องการและออกแบบระบบ & & \cellcolor{gray!45} & \cellcolor{gray!45} & & & & & \\
\hline
พัฒนาและปรับปรุงต้นแบบ Trainy & & & \cellcolor{gray!45} & \cellcolor{gray!45} & \cellcolor{gray!45} & & & \\
\hline
ตรวจเครื่องมือและทดลองนำร่อง & & & & & \cellcolor{gray!45} & \cellcolor{gray!45} & & \\
\hline
เก็บรวบรวมข้อมูลจากผู้ใช้ & & & & & & \cellcolor{gray!45} & \cellcolor{gray!45} & \\
\hline
วิเคราะห์ สรุปผล และจัดทำรายงาน & & & & & & & \cellcolor{gray!45} & \cellcolor{gray!45} \\
\hline
\end{tabularx}
\end{center}

\section{แผนการดำเนินงานและขอบเขตการศึกษา}
\subsection{แผนการดำเนินงาน}
\subsubsection{ศึกษากระบวนการ เอกสาร ผู้ใช้ และปัญหาที่เกิดขึ้นในการจัดการฝึกงาน}
\subsubsection{วิเคราะห์ความต้องการและออกแบบกระบวนการของ Trainy โดยให้ความสำคัญกับบันทึกฝึกงานแบบมีโครงสร้าง}
\subsubsection{พัฒนาหรือปรับปรุงต้นแบบและตรึงรุ่นของระบบสำหรับการทดสอบ}
\subsubsection{ให้ผู้เชี่ยวชาญตรวจเครื่องมืออย่างน้อย 3 คน และทดลองนำร่องกับผู้ใช้ที่ไม่อยู่ในกลุ่มตัวอย่างหลัก}
\subsubsection{ให้ผู้เข้าร่วมทดลองทำภารกิจตามบทบาท บันทึกผลการทำภารกิจ และตอบแบบประเมิน SUS กับคำถามปลายเปิด}
\subsubsection{วิเคราะห์ข้อมูลเชิงพรรณนาและจัดหมวดหมู่ข้อเสนอแนะ แยกผลตามบทบาท}
\subsubsection{สรุปผล ปรับปรุงระบบ และจัดทำรายงานการวิจัย}

\subsection{ขอบเขตการศึกษา}
ศึกษาภายในคณะวิทยาศาสตร์ มหาวิทยาลัยราชภัฏจันทรเกษม ครอบคลุมการสมัครหรือยื่นคำร้อง การอนุมัติและจัดวางนักศึกษา การลงเวลาและคำขอแก้ไข การบันทึกกิจกรรมและความก้าวหน้า การตรวจและรับรองบันทึก การจัดการเอกสาร การแจ้งเตือน การประเมิน และรายงานสรุป ไม่ครอบคลุมการพิสูจน์ผลสัมฤทธิ์ทางการเรียนหรือผลต่อการได้งานในระยะยาว และไม่เปรียบเทียบระหว่างมหาวิทยาลัย

กลุ่มตัวอย่างจำนวน 30 คน ใช้การเลือกแบบเจาะจงร่วมกับโควตาตามบทบาท ได้แก่ นักศึกษา 18 คน อาจารย์นิเทศหรือ\mbox{ผู้ประสานงาน} 5 คน เจ้าหน้าที่ 3 คน และผู้ประกอบการหรือพี่เลี้ยง 4 คน หากเจ้าหน้าที่ที่เกี่ยวข้องจริงมีน้อยกว่า 3 คน จะเชิญทั้งหมดและจัดสรรจำนวนที่เหลือให้กลุ่มผู้ประกอบการ โดยยังคงจำนวนรวม 30 คน

\clearpage
\subsection{ฮาร์ดแวร์ (Hardware) ที่ใช้}
\proposalhardware

\subsection{ซอฟต์แวร์ (Software) ที่ใช้}
\proposalsoftware

\section{ประโยชน์ที่คาดว่าจะได้รับ}
\begin{enumerate}
  \item ได้ข้อกำหนดระบบที่สะท้อนกระบวนการจริงของคณะวิทยาศาสตร์
  \item ได้ระบบกลางที่ช่วยลดการแยกส่วนของข้อมูลและทำให้ติดตามสถานะได้ชัดเจนขึ้น
  \item นักศึกษาสามารถบันทึกหลักฐานการเรียนรู้และความก้าวหน้าได้เป็นระบบ
  \item อาจารย์ เจ้าหน้าที่ และพี่เลี้ยงสามารถตรวจสอบและให้ข้อเสนอแนะผ่านข้อมูลชุดเดียวกัน
  \item ได้หลักฐานด้านการใช้งานและข้อเสนอแนะสำหรับพัฒนาระบบก่อนนำไปใช้จริงในวงกว้างขึ้น
\end{enumerate}

\section{คำนิยามศัพท์เฉพาะ}
\begin{description}[style=nextline,leftmargin=0.8cm,labelindent=0cm,labelwidth=0cm]
  \item[Trainy] เว็บแอปพลิเคชันที่พัฒนาขึ้นเพื่อจัดการกระบวนการฝึกงานและการบันทึกการปฏิบัติงานของนักศึกษา
  \item[บันทึกการฝึกงาน] ข้อมูลที่นักศึกษาบันทึกระหว่างการปฏิบัติงาน เช่น วันที่ เวลา กิจกรรม สิ่งที่เรียนรู้ ปัญหา วิธีแก้ไข ทักษะ และผลลัพธ์
  \item[ผู้ประกอบการหรือพี่เลี้ยง] บุคคลในสถานประกอบการที่กำกับ ตรวจสอบ ให้ข้อเสนอแนะ หรือประเมินนักศึกษาฝึกงาน
  \item[ความสามารถในการใช้งาน] การที่ผู้ใช้สามารถใช้ระบบเพื่อบรรลุภารกิจได้อย่างมีประสิทธิผล มีประสิทธิภาพ และเกิดความพึงพอใจในบริบทที่กำหนด
  \item[กระบวนการตั้งแต่ต้นจนจบ] ขั้นตอนตั้งแต่การเตรียมและสมัครฝึกงานจนถึงการประเมิน สรุปผล และปิดการฝึกงาน
\end{description}

\clearpage
\section{ความคิดเห็นของอาจารย์ที่ปรึกษา}
\noindent\dotfill\\[1.2cm]
\vspace{36pt}
\signatureblock{อาจารย์ที่ปรึกษา}{\proposaladvisor}

\section{ความเห็นของประธานหลักสูตร}
\noindent \fbox{\rule{0pt}{0.55em}\hspace{0.55em}} อนุมัติ\qquad
\fbox{\rule{0pt}{0.55em}\hspace{0.55em}} ไม่อนุมัติ\\[1.2cm]
\vspace{36pt}
\signatureblock{ประธานหลักสูตร}{\pending{\proposalchair}}

\vfill
\noindent\footnotesize\textcolor{gray!70!black}{หมายเหตุ: ข้อความสีแดงและรายละเอียดเทคโนโลยีเป็นรายการที่ต้องตรวจสอบก่อนส่งฉบับจริง}

\end{document}
